% Generated from validated Article Draft Result. Do not hand-edit; revise the structured draft instead.
\section{Model Family Lifecycleが速くなる — Releaseだけでは歴史を追えない}
\noindent\textbf{2025年前半は、model announcementだけを並べると実態を見失う。research preview、API family、retirement announcement、experimental、GAが短い間隔で連続し、availability lifecycleそのものが技術運用上の重要事項になった。}\autocite{src-31799a46d8d0,src-3fa7f4cd510b,src-53f15d483a2d}

2月27日のGPT-4.5はresearch previewとして登場した。これは後続のGPT-4.1とは別のmodel identityであり、pretraining scaleを軸にしたpreviewとして記録する。\autocite{src-3fa7f4cd510b}


4月14日のGPT-4.1 familyはAPI向けにGPT-4.1、mini、nanoを提示し、coding、instruction following、long contextをdeveloper-facingな軸として前面に出した。同日の出来事は単なるGPT-4.5の更新ではない。\autocite{src-31799a46d8d0}


さらに同じ4月14日、GPT-4.5 Preview APIを7月14日にretireする予定が告知された。releaseとdeprecationは異なるobjective eventであり、後者を前者へ畳み込むと利用者が直面したmigration timelineを失う。\autocite{src-31799a46d8d0}


5月22日のClaude Opus 4 / Sonnet 4は新しいmodel generationとして、coding、reasoning、agents、extended thinkingとの接続を前面に出した。Claude 3.7からの流れはfamily continuityを持つ一方、release identityは分けて追う。\autocite{src-4ec03078d316}


6月17日にはGemini 2.5 ProとFlashがGAへ移り、Flash-Liteがpreviewに入った。3月のGemini 2.5 Pro Experimentalと6月のGAは同一系列でも別のavailability eventであり、experimental→GAというservice lifecycleを保存する。\autocite{src-53f15d483a2d}


この半年のmodel familyは「発表日」だけでなく、preview、API-only、GA、retirementを含む時間軸で管理する必要が強まった。能力比較と同時に、採用側にはversion pinningやmigrationを含む運用上の負荷が増えている。\autocite{src-31799a46d8d0,src-53f15d483a2d}


\begin{claimboundary}
本章はrelease identityとavailability chronologyを中心に扱う。各社が掲げる性能指標や優位性は、一次資料の主張と独立再現を区別する。
\end{claimboundary}
